\documentclass[10pt,a4paper]{amsart}
\usepackage[margin=19mm]{geometry}
\usepackage{amsmath,amssymb,mathtools}
\usepackage[scr=rsfs,cal=euler]{mathalfa}
\usepackage{xfrac}
\usepackage[unicode,hidelinks,bookmarks=false]{hyperref}
\newcommand{\ie}{i.e.\ }
\newcommand{\cf}{cf.\ }
\newcommand{\resp}{respectively\ }
\newcommand{\Subsection}{\subsection}
\providecommand{\nobreakdash}{\nobreak\hbox{-}}
\setlength{\emergencystretch}{2em}
\multlinegap0pt
\usepackage{bm}
\usepackage{bbm}
\usepackage{dsfont}
\usepackage{xspace}
\usepackage{cancel}
\usepackage{mathtools}
\usepackage{amsthm}
\makeatletter
\@ifundefined{defn}{\newtheorem{defn}{Defn}}{}
\@ifundefined{prop}{\newtheorem{prop}[defn]{Proposition}}{}
\makeatother
\newcommand{\R}{\mathbb{R}}
\newcommand{\Sp}{{\textsc (S)} }
\newcommand{\di}{{\rm d}}
\title[The local Turnpike Property in Mean Field Control and Games ]{The local Turnpike Property in Mean Field Control and Games with quadratic Hamiltonian}
\author{Marco Cirant, Nicol\`o De Bernardi}
\date{Source: arXiv:2511.18923v2 (CC BY 4.0)}
\begin{document}
\maketitle

\noindent\emph{Source and licence.} The local Turnpike Property in Mean Field Control and Games with quadratic Hamiltonian. Version arXiv:2511.18923v2 (CC BY 4.0), distributed under the Creative Commons Attribution 4.0 International licence, as recorded in the arXiv licence metadata for this exact version and shown on the abstract page https://arxiv.org/abs/2511.18923v2. Full rights notice: licenses/arxiv-2511.18923-CC-BY-4.0.txt.\par\smallskip

\noindent\textit{Editorial scope.} Counted unit: the Prop whose source label is \texttt{phinl1} in the article (source0.tex lines 1274--1288, source label \texttt{phinl1}), delivered together with the article's own complete original proof of it (source0.tex lines 1289--1396, one \texttt{proof} environment of 108 lines). No mathematics is written, repaired or re-derived: statement and proof are the article's own text line for line. Required context retained uncounted: integrab (lines 1246--1252). Every definition, display and label of those lines is retained unchanged. Omitted from this package and not redistributed: 1--1245, 1253--1273, 1397--2642. Nothing inside a retained excerpt is omitted or altered: every retained body is delivered exactly as the article's lines are written, so no omission receipt of any kind accompanies this package. Citations: the retained excerpts carry no citation command at all, so no bibliography entry of the article is transcribed and no reference is redistributed. Reference adaptation: none was needed and none was made; every reference inside the delivered unit resolves inside it, so no unresolved reference or citation is produced. Printed slips kept literal: no printed word, symbol, hypothesis or number of the counted statement or of its proof was altered. Layout and class: the article's own class is \texttt{article} with packages including fontenc, inputenc, babel, amsmath, amssymb, amsthm, calc, accents, bbm, authblk, marginnote, pgfplots, palatino, mathrsfs; the wrapper is the lane's pinned \texttt{amsart} editorial wrapper at 10pt on A4, which restates the theorem-like environments the retained text needs and adds the article's own macro definitions that the retained text needs (\texttt{\textbackslash R}, \texttt{\textbackslash Sp}, \texttt{\textbackslash di}), with extra wrapper packages (bm, bbm, dsfont, xspace, cancel, mathtools). The delivered numbering is the unit's own. Assets: no retained span contains a figure environment, a graphics-inclusion command, a PDF-page inclusion, a table or a TikZ picture; no image file is copied into this package, the package declares no asset dependency and no image file is redistributed with it. Licence: version arXiv:2511.18923v2 of this article is distributed under the Creative Commons Attribution 4.0 International licence, as recorded in the arXiv licence metadata for this exact version and shown on the abstract page https://arxiv.org/abs/2511.18923v2. A full rights notice is delivered at licenses/arxiv-2511.18923-CC-BY-4.0.txt. Characters: every retained excerpt is pure ASCII, so the delivered engine, XeLaTeX with its default Latin Modern text font, reports no missing character.
	\begin{equation}\label{integrab}
	\begin{aligned}
	& \partial_t \mu(t), \mu(t) \in L^1(\Omega), \quad  \mu(t), D \mu(t), D^2 \mu(t), \mu(t) / \bar m \in L^\infty(\Omega), \\
	& \partial_t v(t), v(t), Dv(t), D^2 v(t), D^3 v(t) \in L^\infty(\Omega), \\
	& \int_\Omega \mu(t) = 0
	\end{aligned}
	\end{equation}

	\begin{prop} \label{phinl1}
		Suppose that \eqref{integrab} and \Sp holds. Let  
		$$\tilde{\Phi}(t):=\int_{\Omega}e^{-\delta t} \left( \mu(t,x) v(t,x) + \frac{\mu^2(t,x)}{\bar{m}(x)} \right) \di x.$$
		Then, there exist $\Theta \le 1$ depending on $\bar{m}, \eta$ and a constant $\sigma > 0$ depending on $\eta,C_P,C_f,\bar{m}$ (where $C_P$ is the Poincar\'e constant of $\bar{m}$) such that, if
		$$\sup_{t \in [0,T]}\|\mu(t,\cdot)\|_{L^\infty(\Omega)} + \sup_{t \in [0,T]} \| Dv(t,\cdot) \|_{L^\infty(\Omega)} \le \Theta,$$
		then
		\begin{equation}
			\Phi(T) e^{-(\sigma+\delta)(T-t)} \le \Phi(t) \le \Phi(0)e^{-(\sigma-\delta) t},
			\label{phidelta}
		\end{equation}
		for all $t \in [0,T]$, where
		\begin{equation*}
			\Phi(t):=e^{\delta t} \tilde{\Phi}(t) = \int_{\Omega} \mu(t,x) v(t,x) + \frac{\mu^2(t,x)}{\bar{m}(x)} \, \di x.
		\end{equation*}
	\end{prop}

	\begin{proof}
		We proceed directly with the computation of the derivative of $\tilde{\Phi}$. We first perform some integration by parts, recalling that $D\mu + \mu D\bar{u}= \bar{m}D\left( \frac{\mu}{\bar{m}} \right)$, since $D \bar m = - \bar m D \bar u$. Hypothesis~\eqref{integrab} ensures the operation of integration by parts (in particular when $\Omega = \R^n$) and that derivatives can be moved into integrals.
		\begin{equation*}
			\dot{\tilde{\Phi}}(t)=\int_{\Omega} e^{-\delta t} \left[ (\partial_t \mu) v + \mu (\partial_t v) - \delta \mu v + 2\frac{\mu (\partial_t \mu)}{\bar{m}} - \delta \frac{\mu^2}{\bar{m}} \right] \di x=
		\end{equation*}
		\begin{multline*}
			=\int_{\Omega} e^{-\delta t} \left[ v(\Delta \mu + \text{div}(\mu D\bar{u} + \bar{m}Dv + \mu Dv ) \right] \di x + \\
			+\int_{\Omega} e^{-\delta t} \left[ \mu \left( -\Delta v + \langle Dv,D\bar{u} \rangle - f_m(\bar{m})\mu - \xi \mu^2 + \frac{|Dv|^2}{2} + \delta v \right) \right] - \delta \mu v e^{-\delta t} \di x +  \\
			+ \int_{\Omega} e^{-\delta t} \left[ 2\frac{\mu}{\bar{m}} \left( \text{div}(D\mu + \mu D\bar{u} + \bar{m} Dv + \mu Dv \right) \right] - \delta \frac{\mu^2}{\bar{m}} e^{-\delta t} \di x=
		\end{multline*}
		\begin{multline*}
			=\int_{\Omega} e^{-\delta t} \left[ - \bar{m} |Dv|^2 - \mu |Dv|^2 - f_m(\bar{m})\mu^2 - \xi \mu^3 + \frac{\mu|Dv|^2}{2} - \delta \frac{\mu^2}{\bar{m}} - 2\bar{m} \left| D\left(\frac{\mu}{\bar{m}} \right) \right|^2 \right] \di x +\\
			+\int_{\Omega} e^{-\delta t} \left[- 2 \bar{m} \langle Dv,D\left( \frac{\mu}{\bar{m}} \right) \rangle - 2 \mu \langle Dv,D\left( \frac{\mu}{\bar{m}} \right) \rangle \right] \di x.
		\end{multline*}
		From now on, for the sake of simplicity, we will denote by
		$$z:=\frac{\mu}{\bar{m}}.$$
		Changing all the signs, we get the following equality:
		\begin{multline}
			-\dot{\tilde{\Phi}}(t) = \int_{\Omega} e^{-\delta t} \Big[ \frac{\mu |Dv|^2}{2} + \bar{m}|Dv|^2 +  f_m(\bar{m})\mu^2 + \xi \mu^3 + \delta \frac{\mu^2}{\bar{m}} +\\ + 2\bar{m}|Dz|^2 + 2\bar{m} \langle Dv,Dz \rangle + 2\mu \langle Dv,Dz \rangle \Big] \di x.
			\label{part1}
		\end{multline}
	We now use the hypothesis \Sp with $\mu$. 
%This can be done because
%		\begin{equation}
%			\Delta \mu(t,x) + \text{div}(-\bar{m} Dz(t,x)) + \text{div}(\mu(t,x)D\bar{u}(x)) = 0,
%		\end{equation}
	%	for all $t \in [0,T]$.
		Hence, for all $t \in [0,T]$,
		\begin{equation}
			\int_{\Omega} f_m (\bar{m})\mu^2 + \bar{m}|Dz|^2 + \delta \frac{\mu^2}{\bar{m}} \di x \geq \eta \int_{\Omega} \bar{m} |Dz|^2 \di x,
			\label{coerc}
		\end{equation}
		for some $\eta \in (0,1)$.
		Then, if one plugs~\eqref{coerc} into~\eqref{part1}, this yields
		\begin{equation*}
			\begin{split}
				-\dot{\tilde{\Phi}}(t) &\geq \int_{\Omega} e^{-\delta t} \left[ \frac{\mu |Dv|^2}{2} + \bar{m}|Dv|^2 + \xi \mu^3+ (1+\eta)\bar{m}|Dz|^2  + 2\bar{m} \langle Dv,Dz \rangle + 2\mu \langle Dv,Dz \rangle \right] \di x = \\
				&=\int_{\Omega} e^{-\delta t} \left[ \frac{\mu|Dv|^2}{2}+ \bar{m}|Dv+Dz|^2 + \eta \bar{m}|Dz|^2 +\xi  \mu^3 + 2\mu \langle Dv,Dz \rangle \right] \di x.
			\end{split}
		\end{equation*}
		Since it can be easily shown that (as $\eta \in (0,1)$),
		\begin{equation*}
			\bar{m}|Dv+Dz|^2 + \eta \bar{m} |Dz|^2 \geq \frac{\eta}{4} \bar{m} |Dv|^2 + \frac{\eta}{4} \bar{m} |Dz|^2,
		\end{equation*}
		we can rewrite the whole inequality as follows:
		\begin{equation}
			-\dot{\tilde{\Phi}}(t) \geq \int_{\Omega} e^{-\delta t} \left[ \frac{\eta}{4} \bar{m}|Dv|^2 + \frac{\eta}{4} \bar{m}|Dz|^2 +\xi \mu^3 + 2\mu \langle Dv,Dz \rangle + \frac{\mu|Dv|^2}{2} \right] \di x.
			\label{part2}
		\end{equation}
		Now, we deal with the last three terms of~\eqref{part2} using the smallness hypothesis on $\mu$ and $Dv$ :
		\begin{itemize}
			\item $2 \mu \langle Dv,Dz \rangle \geq - 2 |\mu| |Dv| |Dz| \geq - 2 \Theta |\mu| |Dz| \geq - \Theta \frac{\mu^2}{\bar{m}} - \Theta \bar{m}|Dz|^2;$
			\item $\xi \mu^3 \geq - C_f \left \| \mu(t,\cdot) \right \|_{L^{\infty}(\Omega)} \left \| \bar{m} \right \|_{L^{\infty}(\Omega)} \frac{\mu^2}{\bar{m}} \geq - 
			C_f \|\bar{m}\|_{L^{\infty}(\Omega)} \Theta \frac{\mu^2}{\bar{m}};$
			\item $\frac{1}{2} \mu |Dv|^2 \geq -\frac{1}{2} |\mu| |Dv|^2 \geq - \frac{1}{2} \Theta |\mu| |Dv| \geq - \frac{1}{4} \Theta \frac{\mu^2}{\bar{m}} - \frac{1}{4} \Theta \bar{m} |Dv|^2.$
		\end{itemize}
		We plug all the previous estimates into~\eqref{part2} and we get:
		\begin{multline*}
			-\dot{\tilde{\Phi}}(t) \geq \\
			\int_{\Omega} e^{-\delta t} \bar{m} |Dv|^2 \left[ \frac{\eta}{4} - \frac{\Theta}{4} \right] + e^{-\delta t} \bar{m} |Dz|^2 \left[ \frac{\eta}{4} - \Theta \right]+ e^{-\delta t} \frac{\mu^2}{\bar{m}}  \left[ - \Theta C_f \|\bar{m}\|_{L^{\infty}(\Omega)} - \frac{5 \Theta}{4} \right] \di x \geq
		\end{multline*}
		Now, since $\int_{\R^n} \bar{m}(x) z(t,x) \di x = \int_{\Omega} \mu(t,x) \di x = 0$ for all $t$, by the Poincar\'e inequality we have that
		\[
		 \int_{\Omega} \frac{\mu^2}{\bar{m}} \di x = \int_{\Omega} z^2 \bar{m} \di x \le C_P \int_{\Omega} |Dz|^2 \bar{m} \di x,
		\]
		hence
		\begin{multline*}
			-\dot{\tilde{\Phi}}(t) \geq \\
			\int_{\Omega} e^{-\delta t} \bar{m} |Dv|^2 \left[ \frac{\eta}{4} - \frac{\Theta}{4} \right] + e^{-\delta t} \bar{m} |Dz|^2 \left[ \frac{\eta}{8} - \Theta \right] + e^{-\delta t} \left[\frac{\eta}{8C_P}- \Theta C_f \|\bar{m}\|_{L^{\infty}(\Omega)} -\frac{5\Theta}{4} \right] \frac{\mu^2}{\bar{m}} \di x.
		\end{multline*}
		
		Now, if $$\Theta \le \frac\eta{16} \quad \text{and} \quad \Theta \le \frac{\eta}{16C_P(C_f \|\bar{m}\|_{L^{\infty}(\Omega)} + 5/4)}$$ the following inequality follows:
		\begin{equation}\label{part3}
			\begin{split}
				-\dot{\tilde{\Phi}}(t) & \geq \int_{\Omega} e^{-\delta t}  \left[ \frac{3\eta}{16} \bar{m} |Dv|^2+  \frac\eta{16}  \bar{m} |Dz|^2 +\left(\frac{\eta}{8C_P}- \Theta C_f \|\bar{m}\|_{L^{\infty}(\Omega)} -\frac{5\Theta}{4} \right) \frac{\mu^2}{\bar{m}} \right] \di x \\
				& \geq \int_{\Omega} e^{-\delta t}  \left[ \frac{3\eta}{16} \bar{m} |Dv|^2+  \frac\eta{16}  \bar{m} |Dz|^2+ \frac{\eta}{16C_P} \frac{\mu^2}{\bar{m}} \right] \di x \\
				& \geq \frac{\eta}{32 C_P} \int_{\Omega} e^{-\delta t} \left( \bar{m}\hat{v}^2 + \frac{\mu^2}{\bar{m}} \right) + e^{-\delta t} \frac{\mu^2}{\bar{m}} \di x,
			\end{split}
		\end{equation}
		where $\hat{v}=v-\int_{\Omega} \bar{m}(y) v(t,y) \, \di y$. \\
		Hence, by Young's inequality and the triangle inequality, we get that, for $\sigma = \frac{\eta}{32 C_P}$,
		\begin{equation*}
			-\dot{\tilde{\Phi}}(t) \geq \sigma \int_{\Omega} e^{-\delta t} |\mu \hat{v}| + e^{-\delta t} \left| \frac{\mu^2}{\bar{m}} \right| \di x \geq \sigma e^{-\delta t} \left| \int_{\Omega} \mu \hat{v} + \frac{\mu^2}{\bar{m}} \di x \right| \geq \sigma e^{-\delta t} \left| \int_{\Omega} \mu \hat{v} + \frac{\mu^2}{\bar{m}} \di x \right| = \sigma |\tilde{\Phi}(t)|,
			%\label{inphi}
		\end{equation*}
		since
		\begin{equation*}
			\int_{\Omega} \mu(t,x)\hat{v}(t,x) \di x = \int_{\Omega} \mu v \, \di x - \int_{\Omega} \mu(t,x) \underbrace{\left( \int_{\Omega} \bar{m}(y)v(t,y) \di y \right)}_{=\text{constant w.r.t. } x} \di x = \int_{\Omega} \mu(t,x) v(t,x) \di x.
		\end{equation*}
		Then, we finally obtain that
		\begin{equation*}
			\begin{cases}
				\dot{\tilde{\Phi}}(t) \le \sigma \tilde{\Phi}(t) \\
				\dot{\tilde{\Phi}}(t) \le -\sigma \tilde{\Phi}(t),
			\end{cases}
		\end{equation*}
		from which, after integration on $[t,T]$ and $[0,T]$, one gets
		\begin{equation}\label{tildephiineq}
			\tilde{\Phi}(T) e^{-\sigma(T-t)} \le \tilde{\Phi}(t) \le \tilde{\Phi}(0) e^{-\sigma t},
		\end{equation}
		which implies~\eqref{phidelta}.
		
		Finally, note that integrating \eqref{part3} on $[t_1,t_2]$ gives
		\begin{equation}\label{nuovaint}
		\int_{t_1}^{t_2} \int_{\Omega} e^{-\delta t} \left[\bar{m} |Dv|^2+  \bar{m} |Dz|^2+  \frac{\mu^2}{\bar{m}}\right] \di x \di t \le c [\tilde{\Phi}(t_1)-  \tilde{\Phi}(t_2)],
		\end{equation}
		where $c$ depends on $\eta, C_P$.
	\end{proof}

\end{document}
